\chapter{Architecture technique et flux d'information}
\label{ch:flux}

\section{Introduction}

Les chapitres précédents ont posé les fondations de gouvernance, de sécurité et de stratégie cloud de MABRID. Ce chapitre complète cette vision en décrivant de manière opérationnelle \textbf{comment l'information circule} au sein de la plateforme MarGem : des utilisateurs finaux (acheteurs, vendeurs, visiteurs) jusqu'aux services persistants (base de données, stockage, courriel) en passant par l'application mobile Flutter, l'API FastAPI et le tableau de bord administrateur web séparé.

L'objectif est triple :
\begin{enumerate}
  \item fournir une cartographie lisible pour les parties prenantes non techniques ;
  \item documenter les chemins de données sensibles (authentification, médias, modération) ;
  \item comparer deux modèles de déploiement --- \textbf{cloud Azure} et \textbf{serveur on-premise Ubuntu} --- alignés sur la feuille de route de MABRID.
\end{enumerate}

\section{Vue d'ensemble du système}

\subsection{Acteurs et composants}

MarGem est une plateforme de \emph{découverte et de mise en relation locale} : elle ne traite pas les paiements in-app ; les transactions se concluent hors plateforme. Les composants logiciels principaux sont résumés au tableau~\ref{tab:composants-margem}.

\begin{table}[H]
  \centering
  \caption{Composants logiciels et rôles dans l'écosystème MarGem}
  \label{tab:composants-margem}
  \begin{tabularx}{\textwidth}{@{}lXl@{}}
    \toprule
    \textbf{Composant} & \textbf{Rôle} & \textbf{Technologie} \\
    \midrule
    Application mobile & Interface acheteur / vendeur & Flutter (Android, iOS) \\
    API backend & Logique métier, authentification, modération & FastAPI (Python) \\
    Base de données & Persistance relationnelle & PostgreSQL 16 \\
    Stockage médias & Images produits, avatars & Disque local, Azure Blob ou MinIO \\
    Tableau de bord admin & Gestion staff (hors app mobile) & Web statique (nginx) \\
    Observabilité & Métriques, journaux, alertes & App Insights / Prometheus--Grafana \\
  \bottomrule
  \end{tabularx}
\end{table}

\subsection{Schéma général des flux}

La figure~\ref{fig:flux-overview} synthétise les flux de requêtes entre clients, backend et services externes. Les flèches pleines représentent les échanges de données ; les services en pointillés sont optionnels selon l'environnement de déploiement.

\begin{figure}[H]
  \centering
  \begin{tikzpicture}[
    node distance=1.1cm and 1.4cm,
    comp/.style={rectangle, draw=mabridblue, fill=mabridblue!8, rounded corners=3pt,
      minimum width=2.6cm, minimum height=0.85cm, align=center, font=\small},
    ext/.style={rectangle, draw=mabridaccent, fill=mabridaccent!10, dashed,
      minimum width=2.4cm, minimum height=0.75cm, align=center, font=\footnotesize},
    user/.style={ellipse, draw=mabridblue, fill=white, minimum width=2.2cm, align=center, font=\small},
    arr/.style={-{Stealth[length=2.2mm]}, thick, mabridblue}
  ]
    \node[user] (buyer) {Acheteur\\Vendeur};
    \node[user, right=of buyer] (admin) {Staff\\admin};
    \node[comp, below=1.3cm of buyer] (flutter) {App Flutter\\mobile};
    \node[comp, below=1.3cm of admin] (webadmin) {Admin web\\:8080};
    \node[comp, below=1.6cm of flutter] (api) {API FastAPI\\:8000 / :7226};
    \node[comp, below left=1.2cm and 0.3cm of api] (pg) {PostgreSQL};
    \node[comp, below right=1.2cm and 0.3cm of api] (media) {Stockage\\médias};
    \node[ext, right=2.2cm of api] (smtp) {SMTP};
    \node[ext, above=0.5cm of smtp] (maps) {Google Maps};
    \node[ext, below=0.5cm of smtp] (stripe) {Stripe};

    \draw[arr] (buyer) -- (flutter);
    \draw[arr] (admin) -- (webadmin);
    \draw[arr] (flutter) -- node[right, font=\footnotesize] {HTTPS + JWT} (api);
    \draw[arr] (webadmin) -- node[above, font=\footnotesize] {HTTPS + JWT} (api);
    \draw[arr] (api) -- (pg);
    \draw[arr] (api) -- (media);
    \draw[arr, dashed] (api) -- (smtp);
    \draw[arr, dashed] (flutter) -- (maps);
    \draw[arr, dashed] (api) -- (stripe);
  \end{tikzpicture}
  \caption{Vue d'ensemble des flux d'information MarGem}
  \label{fig:flux-overview}
\end{figure}

\section{Architectures de déploiement}

\subsection{Déploiement cloud Microsoft Azure}

La figure~\ref{fig:arch-azure} présente l'architecture cible de production sur Azure. Elle s'inscrit dans une \emph{zone d'atterrissage} (\textit{landing zone}) structurée autour d'un abonnement dédié (\texttt{rg-margem-prod}), d'un réseau virtuel isolé et de points de terminaison privés pour les données sensibles.

\begin{figure}[H]
  \centering
  \IfFileExists{figures/architecture-azure-prod.png}{%
    \includegraphics[width=\textwidth]{architecture-azure-prod.png}%
  }{%
    \fcolorbox{mabridblue!40}{mabridblue!5}{\parbox{0.9\textwidth}{\centering\vspace{2cm}\textbf{Figure Azure}\\\texttt{figures/architecture-azure-prod.png}\vspace{2cm}}}%
  }
  \caption{Architecture de déploiement MarGem sur Microsoft Azure (production)}
  \label{fig:arch-azure}
\end{figure}

Les éléments structurants sont les suivants :
\begin{itemize}
  \item \textbf{Réseau} --- \texttt{margem-vnet}, groupes de sécurité réseau (NSG), points de terminaison privés vers PostgreSQL, Key Vault et Blob Storage.
  \item \textbf{Calcul} --- Azure Container Apps (\texttt{margem-api}, FastAPI, port~8000) avec identité managée et montée en charge automatique.
  \item \textbf{Données} --- PostgreSQL Flexible Server v16 (\texttt{margem}) et conteneur Blob \texttt{margem-media} (redondance locale LRS).
  \item \textbf{Identité} --- Azure AD, \gls{rbac}, journaux d'audit ; accès moindre privilège via identité utilisateur assignée.
  \item \textbf{Observabilité} --- Application Insights et Log Analytics pour métriques, exceptions et alertes.
  \item \textbf{Intégrations} --- Stripe (abonnements), fournisseur SMTP (courriels transactionnels), Google Maps (géolocalisation optionnelle).
\end{itemize}

Le client mobile et le tableau de bord admin (optionnel) communiquent avec l'API via \textbf{HTTPS REST} ; les secrets d'exécution sont injectés depuis Azure Key Vault (\texttt{database-url}, \texttt{jwt-secret}, \texttt{storage-conn}, etc.).

\subsection{Déploiement on-premise (Ubuntu)}

Pour les phases pilotes, les démonstrations académiques ou les déploiements à coût maîtrisé, MarGem peut être hébergé sur un \textbf{serveur Ubuntu local} (réseau domestique ou centre de données interne). La figure~\ref{fig:arch-onprem} illustre cette variante.

\begin{figure}[H]
  \centering
  \IfFileExists{figures/architecture-onprem-ubuntu.png}{%
    \includegraphics[width=\textwidth]{architecture-onprem-ubuntu.png}%
  }{%
    \fcolorbox{mabridblue!40}{mabridblue!5}{\parbox{0.9\textwidth}{\centering\vspace{2cm}\textbf{Figure On-Prem}\\\texttt{figures/architecture-onprem-ubuntu.png}\vspace{2cm}}}%
  }
  \caption{Architecture de déploiement MarGem sur serveur Ubuntu (on-premise)}
  \label{fig:arch-onprem}
\end{figure}

Comparée au cloud Azure, cette architecture repose sur :
\begin{itemize}
  \item \textbf{Périmètre réseau local} --- pare-feu UFW/iptables, proxy inverse Nginx ou Traefik, certificats Let's Encrypt, Fail2Ban.
  \item \textbf{Orchestration} --- Docker Compose (développement et serveur domestique) ou Kubernetes on-premise à l'échelle.
  \item \textbf{Secrets} --- coffre HashiCorp Vault ou variables chiffrées (\texttt{.env.home} hors dépôt).
  \item \textbf{Stockage} --- MinIO (compatible S3) ou système de fichiers local (\texttt{/data/media}).
  \item \textbf{Observabilité} --- pile Prometheus, Grafana, ELK/Loki en lieu et place d'Application Insights.
\end{itemize}

Le choix entre Azure et on-premise n'est pas exclusif : la feuille de route prévoit un \textbf{serveur domestique} pour le prototypage (API sur le port~7226, admin sur~8080) et une migration progressive vers Azure pour la montée en charge nationale.

\section{Flux d'authentification et de gestion de session}

\subsection{Séquence d'inscription et de connexion}

L'authentification repose sur des jetons \gls{jwt} (accès court + rafraîchissement rotatif). La figure~\ref{fig:flux-auth} décrit les échanges principaux entre l'application mobile et l'API.

\begin{figure}[H]
  \centering
  \begin{tikzpicture}[
    x=0.42cm, y=0.65cm,
    actor/.style={rectangle, draw=mabridblue, minimum width=2cm, minimum height=0.7cm, font=\footnotesize},
    msg/.style={font=\scriptsize},
    lifeline/.style={dashed, mabridblue!50}
  ]
    \node[actor] (app) at (0,0) {App Flutter};
    \node[actor] (api) at (14,0) {API FastAPI};
    \node[actor] (db) at (28,0) {PostgreSQL};
    \draw[lifeline] (0,-0.5) -- (0,-16);
    \draw[lifeline] (14,-0.5) -- (14,-16);
    \draw[lifeline] (28,-0.5) -- (28,-16);

    \draw[-{Stealth}] (0,-1.5) -- node[msg, above, sloped] {POST /auth/register} (14,-1.5);
    \draw[-{Stealth}] (14,-2.5) -- node[msg, above, sloped] {INSERT user (rôle=acheteur)} (28,-2.5);
    \draw[-{Stealth}] (28,-3.5) -- (14,-3.5);
    \draw[-{Stealth}] (14,-4.5) -- node[msg, above, sloped] {access + refresh JWT} (0,-4.5);

    \draw[-{Stealth}] (0,-6.5) -- node[msg, above, sloped] {POST /auth/login} (14,-6.5);
    \draw[-{Stealth}] (14,-7.5) -- node[msg, above, sloped] {Vérifier mot de passe} (28,-7.5);
    \draw[-{Stealth}] (28,-8.5) -- (14,-8.5);
    \draw[-{Stealth}] (14,-9.5) -- node[msg, above, sloped] {JWT} (0,-9.5);

    \draw[-{Stealth}] (0,-11.5) -- node[msg, above, sloped] {Requêtes API (Bearer)} (14,-11.5);
    \draw[-{Stealth}] (14,-12.5) -- node[msg, above, sloped] {Lecture / écriture} (28,-12.5);
    \draw[-{Stealth}] (28,-13.5) -- (14,-13.5);
    \draw[-{Stealth}] (14,-14.5) -- node[msg, above, sloped] {JSON} (0,-14.5);
  \end{tikzpicture}
  \caption{Flux d'authentification : inscription, connexion et requêtes authentifiées}
  \label{fig:flux-auth}
\end{figure}

\subsection{Contrôles de sécurité associés}

Au-delà du schéma séquentiel, plusieurs mécanismes renforcent la session (détaillés au chapitre~\ref{ch:app-security}) :
\begin{itemize}
  \item hachage bcrypt, verrouillage après échecs répétés, détection de réutilisation de refresh token ;
  \item \gls{mfa} obligatoire pour les comptes staff en production ;
  \item révocation globale via version de jeton (\texttt{tv}) lors d'un changement de mot de passe.
\end{itemize}

\section{Flux fonctionnels de l'application mobile}

\subsection{Parcours acheteur et vendeur}

La figure~\ref{fig:flux-mobile} regroupe les principaux parcours côté client mobile. L'application ne propose \textbf{pas} d'interface d'administration : la modération et la configuration des places de marché sont réservées au tableau de bord web.

\begin{figure}[H]
  \centering
  \begin{tikzpicture}[
    node distance=0.55cm and 0.9cm,
    screen/.style={rectangle, draw=mabridblue, fill=mabridblue!6, rounded corners=2pt,
      minimum width=2.1cm, minimum height=0.65cm, font=\scriptsize, align=center},
    api/.style={rectangle, draw=mabridaccent, fill=mabridaccent!8, minimum width=2.4cm,
      minimum height=0.65cm, font=\scriptsize, align=center},
    arr/.style={-{Stealth[length=1.8mm]}, mabridblue, thick}
  ]
    \node[screen] (home) {Accueil\\Découverte};
    \node[screen, right=of home] (search) {Recherche\\Carte};
    \node[screen, right=of search] (prod) {Fiche\\produit};
    \node[screen, below=1.1cm of home] (msg) {Messagerie};
    \node[screen, below=1.1cm of search] (sell) {Tableau de\\bord vendeur};
    \node[screen, below=1.1cm of prod] (prof) {Profil};

    \node[api, below=2.4cm of search] (disc) {/sellers, /search\\/favorites};
    \node[api, left=0.6cm of disc] (chat) {/messages/*};
    \node[api, right=0.6cm of disc] (seller) {/seller/*, /uploads};

    \draw[arr] (home) -- (disc);
    \draw[arr] (search) -- (disc);
    \draw[arr] (prod) -- (disc);
    \draw[arr] (msg) -- (chat);
    \draw[arr] (sell) -- (seller);
    \draw[arr] (prof) -- (disc);
  \end{tikzpicture}
  \caption{Correspondance écrans mobiles $\rightarrow$ endpoints API}
  \label{fig:flux-mobile}
\end{figure}

\begin{table}[H]
  \centering
  \caption{Exemples d'actions utilisateur et persistance des données}
  \label{tab:flux-actions}
  \begin{tabularx}{\textwidth}{@{}lXX@{}}
    \toprule
    \textbf{Action} & \textbf{Endpoint API} & \textbf{Persistance} \\
    \midrule
    Parcourir les commerces & \texttt{GET /sellers}, \texttt{/sellers/map} & PostgreSQL \\
    Rechercher & \texttt{GET /search} & PostgreSQL \\
    Ajouter aux favoris & \texttt{POST /favorites} & PostgreSQL \\
    Contacter un vendeur & \texttt{POST /messages/...} & PostgreSQL \\
    Publier une annonce & \texttt{POST /seller/products} & PostgreSQL + médias \\
    Téléverser une image & \texttt{POST /uploads} & Disque / Blob \\
    \bottomrule
  \end{tabularx}
\end{table}

\section{Tableau de bord administrateur web (séparé)}

\subsection{Séparation des préoccupations}

Conformément au principe de moindre privilège et à la réduction de surface d'attaque, l'administration staff n'est \textbf{pas intégrée} à l'application mobile Flutter. Un conteneur nginx dédié sert l'interface sur le port~\textbf{8080} (configurable), distinct de l'API (port~8000 ou~7226 en serveur domestique).

\subsection{Flux et couches de sécurité}

La figure~\ref{fig:flux-admin} et le tableau~\ref{tab:admin-security} décrivent le chemin d'accès depuis un navigateur mobile ou ordinateur portable sur le réseau local.

\begin{figure}[H]
  \centering
  \begin{tikzpicture}[
    node distance=0.7cm and 1.2cm,
    box/.style={rectangle, draw=mabridblue, fill=mabridblue!7, rounded corners=2pt,
      minimum width=2.5cm, minimum height=0.7cm, align=center, font=\small},
    arr/.style={-{Stealth}, thick, mabridblue}
  ]
    \node[box] (browser) {Navigateur\\:8080};
    \node[box, right=of browser] (nginx) {nginx\\Basic Auth};
    \node[box, right=of nginx] (spa) {Admin SPA\\HTML/JS};
    \node[box, below=1.4cm of nginx] (api) {API :7226\\/admin/*};
    \node[box, below=1.4cm of api] (db) {PostgreSQL};

    \draw[arr] (browser) -- node[above, font=\footnotesize] {HTTPS/LAN} (nginx);
    \draw[arr] (nginx) -- (spa);
    \draw[arr] (spa) -- node[right, font=\footnotesize] {JWT + Origin} (api);
    \draw[arr] (api) -- (db);
  \end{tikzpicture}
  \caption{Flux du tableau de bord administrateur séparé}
  \label{fig:flux-admin}
\end{figure}

\begin{table}[H]
  \centering
  \caption{Couches de sécurité du tableau de bord web}
  \label{tab:admin-security}
  \begin{tabularx}{\textwidth}{@{}lX@{}}
    \toprule
    \textbf{Couche} & \textbf{Contrôle} \\
    \midrule
    Réseau & Liste blanche IP privées (192.168/10/172.16) ; pas d'exposition publique \\
    nginx & Authentification HTTP Basic (mot de passe de porte) \\
    En-têtes & CSP stricte, \texttt{X-Frame-Options: DENY}, limitation de débit \\
    API & Garde d'origine (\texttt{Origin}), rôles \texttt{admin}/\texttt{support}, \gls{mfa} \\
    Session & Jeton en \texttt{sessionStorage} (effacé à la fermeture de l'onglet) \\
    API embarquée & Désactivée sur serveur domestique (\texttt{/admin} non servi sur le port API) \\
    \bottomrule
  \end{tabularx}
\end{table}

\section{Flux de téléversement des médias}

Les images de produits et d'avatars suivent un flux en plusieurs étapes afin d'éviter l'upload direct non contrôlé vers le stockage (figure~\ref{fig:flux-upload}).

\begin{figure}[H]
  \centering
  \begin{tikzpicture}[
    node distance=1.5cm,
    box/.style={rectangle, draw=mabridblue, fill=mabridblue!8, minimum width=2.2cm,
      minimum height=0.75cm, align=center, font=\small},
    arr/.style={-{Stealth}, thick, mabridblue}
  ]
    \node[box] (app) {App\\vendeur};
    \node[box, right=of app] (api) {API};
    \node[box, right=of api] (store) {Stockage\\local/Blob};
    \node[box, right=of store] (db) {PostgreSQL};

    \draw[arr] (app) -- node[above, font=\footnotesize] {1. Token upload} (api);
    \draw[arr] (app) to[bend left=15] node[above, font=\footnotesize] {3. POST /uploads} (api);
    \draw[arr] (api) -- node[above, font=\footnotesize] {4. Écriture} (store);
    \draw[arr] (api) -- node[above, font=\footnotesize] {6. URL média} (db);
    \draw[arr] (api) to[bend left=20] node[below, font=\footnotesize] {5. URL publique} (app);
  \end{tikzpicture}
  \caption{Flux de téléversement et persistance des médias}
  \label{fig:flux-upload}
\end{figure}

\section{Déploiement serveur domestique (cas d'étude)}

Le déploiement sur ordinateur portable Linux (\texttt{docker-compose.home.yml}) illustre un scénario pédagogique et opérationnel fréquent au Maroc : API accessible sur le réseau Wi-Fi local, sans dépendance cloud obligatoire.

\begin{figure}[H]
  \centering
  \begin{tikzpicture}[
    node distance=1.2cm and 1.5cm,
    device/.style={rectangle, draw=mabridblue, rounded corners=3pt, fill=white,
      minimum width=2.4cm, minimum height=0.8cm, align=center, font=\small},
    docker/.style={rectangle, draw=mabridaccent, thick, fill=mabridaccent!5,
      minimum width=7cm, minimum height=3.2cm},
    arr/.style={-{Stealth}, thick}
  ]
    \node[device] (phone) {Téléphone\\Wi-Fi};
    \node[docker, below=1.5cm of phone, xshift=1.5cm] (dock) {};
    \node[font=\small\bfseries, mabridblue] at ([yshift=1.1cm]dock.center) {Docker (piocco)};
    \node[device] at ([xshift=-2cm, yshift=-0.2cm]dock.center) (api) {API\\:7226};
    \node[device] at ([xshift=0cm, yshift=-0.2cm]dock.center) (adm) {Admin\\:8080};
    \node[device] at ([xshift=2cm, yshift=-0.2cm]dock.center) (pg) {Postgres};

    \draw[arr, mabridblue] (phone) -- node[left, font=\footnotesize] {App} (api);
    \draw[arr, mabridblue] (phone) to[bend left=25] node[right, font=\footnotesize] {Admin web} (adm);
    \draw[arr, mabridaccent] (adm) -- (api);
    \draw[arr, mabridaccent] (api) -- (pg);
  \end{tikzpicture}
  \caption{Réseau domestique : application mobile et admin web vers le serveur local}
  \label{fig:flux-home}
\end{figure}

\textbf{Recommandations opérationnelles} :
\begin{itemize}
  \item ne pas exposer le port~8080 via tunnel Cloudflare ; réserver le tunnel à l'API mobile si nécessaire ;
  \item configurer \texttt{CORS\_ORIGINS} en liste séparée par virgules (sans crochets JSON) ;
  \item promouvoir les comptes staff via SQL ou script dédié (\texttt{promote-admin}) ;
  \item activer \texttt{ADMIN\_BASIC\_AUTH\_*} et \texttt{ADMIN\_REQUIRE\_STAFF\_MFA=true} dans \texttt{.env.home}.
\end{itemize}

\section{Comparatif des modèles de déploiement}

\begin{table}[H]
  \centering
  \caption{Comparaison Azure production vs serveur domestique Ubuntu}
  \label{tab:compare-deploy}
  \begin{tabularx}{\textwidth}{@{}lXX@{}}
    \toprule
    \textbf{Critère} & \textbf{Azure (fig.~\ref{fig:arch-azure})} & \textbf{On-prem / domicile (fig.~\ref{fig:arch-onprem})} \\
    \midrule
    Mise à l'échelle & Auto-scaling Container Apps & Manuelle / verticale \\
    Secrets & Key Vault + RBAC & Vault local ou \texttt{.env.home} \\
    Stockage médias & Blob Storage LRS & Disque local ou MinIO \\
    Observabilité & App Insights, Log Analytics & Prometheus, Grafana, ELK \\
    Coût initial & Abonnement cloud & Matériel existant \\
    Sécurité périmétrique & NSG, Private Endpoints & UFW, nginx, IP LAN \\
    Cas d'usage & Production nationale & Prototype, démo, LAN \\
    \bottomrule
  \end{tabularx}
\end{table}

\section{Conclusion du chapitre}

Ce chapitre a cartographié les flux d'information de MarGem de bout en bout : clients mobiles, API, persistance, stockage, intégrations tierces et tableau de bord administrateur séparé. Les figures~\ref{fig:arch-azure} et~\ref{fig:arch-onprem} ancrent ces flux dans deux réalités d'hébergement complémentaires --- cloud managé et serveur Ubuntu --- tandis que les schémas séquentiels et fonctionnels facilitent l'audit, l'onboarding des développeurs et la communication avec les encadrants académiques.

La cohérence avec les chapitres de sécurité (ch.~\ref{ch:security},~\ref{ch:app-security}) et de gouvernance cloud (ch.~\ref{ch:cloud}) garantit que chaque flèche du diagramme correspond à un contrôle documenté : authentification forte, séparation admin/mobile, chiffrement des secrets et traçabilité des opérations sensibles.
